\documentclass{article}

\usepackage[T1]{fontenc}
\usepackage{amsmath,amssymb}
\usepackage{xurl}
\usepackage[hidelinks]{hyperref}
\setlength{\emergencystretch}{2em}

% Shared commands for notes in docs/paper-gaps/.

\DeclareUnicodeCharacter{2080}{\ifmmode{}_{0}\else\textsubscript{0}\fi}
\DeclareUnicodeCharacter{2081}{\ifmmode{}_{1}\else\textsubscript{1}\fi}
\DeclareUnicodeCharacter{2082}{\ifmmode{}_{2}\else\textsubscript{2}\fi}
\DeclareUnicodeCharacter{2099}{\ifmmode{}_{n}\else\textsubscript{n}\fi}

\newcommand{\Lean}{\textsc{Lean}}
\ExplSyntaxOn
\NewDocumentCommand{\leanid}{m}{
  \ifmmode
    \textsf{\detokenize{#1}}
  \else
    \group_begin:
    \urlstyle{sf}
    \tl_set:Nn \l_tmpa_tl {#1}
    \tl_replace_all:Nnn \l_tmpa_tl {\_} {_}
    \exp_args:NV \path \l_tmpa_tl
    \group_end:
  \fi
}
\ExplSyntaxOff
\providecommand{\tr}{\operatorname{tr}}
\newcommand{\tnleanIssueBase}{https://github.com/LionSR/TNLean/issues/}
\newcommand{\tnleanPrBase}{https://github.com/LionSR/TNLean/pull/}
\newcommand{\tnleanDocBase}{https://sirui-lu.com/TNLean/docs/}
\newcommand{\ghissue}[1]{\href{\tnleanIssueBase#1}{GitHub issue \##1}}
\newcommand{\ghpr}[1]{\href{\tnleanPrBase#1}{PR \##1}}
\newcommand{\doclink}[1]{\href{\tnleanDocBase#1.html}{\path{#1}}}

% Dirac notation and the identity operator, as in blueprint/src/macros/common.tex.
% \providecommand keeps note-local definitions authoritative.
\usepackage{dsfont}
\providecommand{\ket}[1]{|#1\rangle}
\providecommand{\bra}[1]{\langle #1|}
\providecommand{\braket}[2]{\langle #1 | #2 \rangle}
\providecommand{\ketbra}[2]{|#1\rangle\!\langle #2|}
\providecommand{\Id}{\mathds{1}}
\providecommand{\one}{\mathds{1}}
\providecommand{\C}{\mathbb{C}}
\providecommand{\MN}[1]{M_{#1}(\C)}
\providecommand{\MD}{M_D(\C)}

% External referencing for paper sources.  The build helper writes sanitized
% auxiliary files under build/paper-aux/ and excludes duplicated source labels.
\usepackage{xr}

% Import a generated paper auxiliary file with a caller-supplied label prefix.
% The candidate paths support builds from docs/paper-gaps/, the repository root,
% and build/paper-aux/ itself.
\newcommand{\importpaperaux}[2]{%
  \IfFileExists{../../build/paper-aux/#2.aux}{%
    \externaldocument[#1]{../../build/paper-aux/#2}%
  }{%
    \IfFileExists{build/paper-aux/#2.aux}{%
      \externaldocument[#1]{build/paper-aux/#2}%
    }{%
      \IfFileExists{#2.aux}{%
        \externaldocument[#1]{#2}%
      }{}%
    }%
  }%
}

% General external reference macros
\newcommand{\paperref}[2]{\ref{#1-#2}}
\newcommand{\papereqref}[2]{(\ref{#1-#2})}

% CPSV16 source (1606.00608 / MPDO-22-12-17-2.tex) external document and shortcuts
\importpaperaux{cpsv16-}{1606.00608/MPDO-22-12-17-2-tnlean}
\newcommand{\cpsvref}[1]{\paperref{cpsv16}{#1}}
\newcommand{\cpsveqref}[1]{\papereqref{cpsv16}{#1}}

% Verdict marker: \gapnote{<kind>}{<status>}. Kind is the policy
% classification (clarification, local-correction, scope-restriction,
% unfaithful, false-source, open-gap); status is open, wip (elimination
% underway), resolved, or historical. Machine-read by the site index;
% prints a verdict line.
\newcommand{\gapnote}[2]{\par\noindent\textbf{Verdict:} #1 (#2).\par}


\title{The Sector Decomposition in the GSNNCH Definition of
Cirac--P\'erez-Garc\'ia--Schuch--Verstraete}
\author{}
\date{2026-07-15}

\begin{document}
\maketitle
\gapnote{local-correction}{resolved}

\section{Source definition}

Definition~4.8 of Ref.~\cite{Cirac2016MPDO_arXiv} defines a density operator
on a periodic chain to be a Gibbs state of a nearest-neighbor commuting
Hamiltonian when
\begin{align}
    \rho^{(N)}
        &\propto
        \bigoplus_x n_x
        \exp\left(-\sum_{j=1}^N\tau_j(h^{(x)})\right),
    \label{eq:cpsv16_gsnnch_sector_decomposition_source_exponential_form}\\
    n_x
        &\in\mathbb N,
    \label{eq:cpsv16_gsnnch_sector_decomposition_source_natural_multiplicity}\\
    [h^{(x)},\tau_1(h^{(x)})]
        &=0.
    \label{eq:cpsv16_gsnnch_sector_decomposition_source_local_commutation}
\end{align}
The one-site spaces associated with distinct outer labels $x$ are orthogonal;
see Definition~4.8, lines~829--842 of
Ref.~\cite{Cirac2016MPDO_arXiv}.

Using the projector-limit convention stated immediately after that definition,
the same class can be written with positive semidefinite two-site bonds
$B^{(x)}$:
\begin{align}
    \rho^{(N)}
        &\propto
        \bigoplus_x n_x
        \prod_{j=1}^N\tau_j(B^{(x)}),
    \label{eq:cpsv16_gsnnch_sector_decomposition_source_bond_product_form}\\
    B^{(x)}
        &\geq0.
    \label{eq:cpsv16_gsnnch_sector_decomposition_source_positive_bond}
\end{align}
The neighboring translates commute; see lines~843--850 of
Ref.~\cite{Cirac2016MPDO_arXiv}. The outer direct sum and the natural
multiplicities are part of the source definition. Replacing the exponential
presentation by positive semidefinite bonds does not remove them.

\section{Present formal boundary}

The formal sector decomposition records a finite outer sector type, orthogonal
one-site sector projections, the natural multiplicities, and one positive
commuting bond in every sector. The finite-chain and tensor-level predicates
therefore retain the direct sum and multiplicities of Definition~4.8 in
Ref.~\cite{Cirac2016MPDO_arXiv}. The finite-chain predicate requires positivity
and unit trace. The tensor-level predicate applies it to the normalized
operator at every length $N\geq2$.\footnote{\raggedright Formal structure:
\leanid{MPOTensor.GSNNCHData}. Formal declarations:
\leanid{MPOTensor.IsGSNNCHAt}, \leanid{MPOTensor.IsGSNNCH}, and
\leanid{MPOTensor.normalizedMPO}.}

The auxiliary single-bond presentation is stored separately.\footnote{\raggedright
Formal structure: \leanid{MPOTensor.CommutingBondProductData}.} It has the
form
\begin{align}
    \rho^{(N)}
        &=c_N\prod_{j=1}^N\tau_j(B_N),
    \label{eq:cpsv16_gsnnch_sector_decomposition_single_bond_presentation}\\
    c_N
        &>0.
    \label{eq:cpsv16_gsnnch_sector_decomposition_single_bond_positive_scalar}
\end{align}
A theorem embeds this presentation as the one-sector case of the source
definition. No converse theorem is claimed that extracts a single bond from
arbitrary source sector data.

The represented finite-chain sector sum has the asserted analytic properties:
all cyclic translates of each sector bond commute pairwise, the represented
operator is positive semidefinite and cyclic-shift invariant, and a tensor
whose finite-chain operators have the sector form with nonzero traces satisfies
the full density-operator predicate. In particular, every injective tensor
satisfying the saturated area law is proved GSNNCH through a one-sector
witness.\footnote{\raggedright Formal declarations:
\leanid{MPOTensor.GSNNCHData.bondAt_comm},
\leanid{MPOTensor.GSNNCHData.unnormalizedState_posSemidef},
\leanid{MPOTensor.GSNNCHData.unnormalizedState_submatrix_rotateConfig},
\leanid{MPOTensor.isGSNNCH_of_hasGSNNCHForm}, and
\leanid{MPOTensor.isGSNNCH_of_isInjective_of_isSAL}.}
This one-sector form is the commuting-product conclusion of
Proposition~C.8 in Ref.~\cite{Cirac2016MPDO_arXiv} for an injective tensor.

\section{Elimination plan}

The first part of the original plan is complete: positivity and translation
invariance of the represented finite-chain sector sum follow from the stored
sector data. Appendix~C.2 of Ref.~\cite{Cirac2016MPDO_arXiv} has also been
formalized through the BNT projector-selection formula. The natural
multiplicities are retained; positive-length compression proves that every
absorbed BNT representative generates MPDOs; the direct-sum entropy identity
proves SAL for each representative; and the canonical block equation proves
ZCL for each representative.\footnote{\raggedright Formal declarations:
\leanid{MPOTensor.exists_bntProjectorSelection_positiveLength_of_sameMPV₂Pos_isSAL},
\leanid{MPOTensor.commonWeightAbsorbedBasisMPOTensor_isMPDO_of_sameMPV₂Pos_isSAL},
\leanid{MPOTensor.commonWeightAbsorbedBasisMPOTensor_isSAL_of_sameMPV₂Pos}, and
\leanid{MPOTensor.commonWeightAbsorbedBasisMPOTensor_isSourceZCL}.}

The general outer-sector construction is also complete. Pairwise orthogonal
one-site projections, supported positive commuting bonds, and exact sector
product identities produce the source GSNNCH form with the prescribed natural
multiplicities.\footnote{\raggedright Formal declaration:
\leanid{MPOTensor.hasGSNNCHForm_of_orthogonalCommutingSectorFamily}.}

The support-preserving sector argument is complete as well. Every absorbed BNT
representative has an exact isometric restriction to the range of its matching
projector. The restricted tensor is injective, SAL passes from the ambient
representative to the restriction through the sitewise support isometry, and
the positive Proposition~C.8 bond of
Ref.~\cite{Cirac2016MPDO_arXiv} lifts to a positive two-site operator supported
on the tensor square of the projector. Its cyclic translates commute pairwise
on every periodic chain of length at least two, and their full periodic product
equals the ambient finite-chain operator exactly, with normalization scalar
one as in the restricted product.\footnote{\raggedright Formal declarations:
\leanid{MPOTensor.PhysicalSupportRestrictionData},
\leanid{MPOTensor.nonempty_physicalSupportRestrictionData_commonWeightAbsorbedBasis},
\leanid{MPOTensor.PhysicalSupportRestrictionData.restricted_isSAL}, and
\leanid{MPOTensor.PhysicalSupportRestrictionData.exists_etaLocalStructureData_lifted_supported}.}

Choosing supported products for all absorbed BNT representatives gives an
orthogonal commuting sector family once the required one-site supports are
available. A source-context corollary derives pairwise two-sided supports from
common copy weights, one-site spanning, positive neighboring-operator
factorizations, and the unstarred layer identity. For distinct sectors $x$ and
$y$, that local identity is
\begin{align}
    \mathcal{K}_y[\beta_y,\alpha_y]
    \mathcal{K}_x[\beta_x,\alpha_x]
        &=0
        \qquad (x\neq y).
    \label{eq:cpsv16_gsnnch_sector_decomposition_unstarred_layer_annihilation}
\end{align}
This is an ordered product, not an adjoint product.

Under the standing biCF hypothesis, each BNT tensor is injective. The source
identities \texttt{AppUkU=rl} and \texttt{Appetakhetc} in Appendix~C.2 of
Ref.~\cite{Cirac2016MPDO_arXiv} supply positive neighboring operators and
therefore an MPDO in each sector. Hermiticity identifies every sector tensor
with its physical adjoint through the injective fundamental theorem. The span
of the physical slices is consequently closed under conjugate transpose. The
unstarred ordered annihilation identities then make the slice-column spaces
orthogonal, and their support projections fix every sector slice on both
sides.\footnote{\raggedright Formal declarations:
\leanid{MPOTensor.IsBNTLayerOrthogonal},
\leanid{MPOTensor.PhysicalSectorFactorization.isMPDO_of_neighboringOperator_pos},
\leanid{MPOTensor.exists_physicalSlice_conjTranspose_eq_sum_of_isInjective_isMPDO},
\leanid{MPOTensor.exists_pairwise_orthogonal_twoSided_physicalSupport}, and
\leanid{MPOTensor.exists_pairwise_orthogonal_twoSided_physicalSupport_commonWeightAbsorbedBasis}.}

This result formalizes the pairwise support consequence of the earlier
Case~II data. The ambient factorization \texttt{AppUkU=rl}, together with
\texttt{Appetakhetc}, supplies the canonical injective active restriction and
its positive physical-sector factorization. Its positive commuting bond lifts
to a bond $B_x$ on the original one-site space satisfying
\begin{align}
    (P_x\otimes P_x)B_x(P_x\otimes P_x)
        &=B_x.
    \label{eq:cpsv16_gsnnch_sector_decomposition_supported_lifted_bond}
\end{align}
All periodic translates of $B_x$ commute, and their product equals
$\rho^{(N)}(\mathcal K_x)$ exactly for every $N\geq2$. Applying the construction
to pairwise orthogonal projections gives the orthogonal commuting sector family
required for the outer GSNNCH sum.\footnote{\raggedright Formal declarations:
\leanid{MPOTensor.PhysicalSectorFactorization.exists_etaLocalStructureData_supported_of_twoSidedPhysicalSlice}
and
\leanid{MPOTensor.nonempty_orthogonalCommutingSectorFamily_of_ambientPhysicalSectorFactorization}.}

The active factor supports in the ambient factorization
\texttt{AppUkU=rl} have been assembled into the canonical physical
restriction. Each active sector retains its left and right joint column
supports, while an inactive sector contributes a zero-dimensional coordinate
space. The direct-sum inclusion has orthonormal columns and range exactly
$P_x$. Compression has the corresponding sectorwise left--right factorization,
and re-embedding recovers $\mathcal K_x$ exactly. No identity factor in an
unused direction or complementary zero summand is introduced. This statement
includes the case of an empty active family.

Every neighboring operator of the compressed factorization is identified with
the explicit congruence of the corresponding ambient operator in
\texttt{Appetakhetc}. Positive semidefiniteness follows from congruence.
Therefore the canonical restriction has a positive physical-sector
factorization indexed precisely by the active ambient sectors, including the
vacuous empty-family case.

The canonical joint physical support is contained in every orthogonal
projection satisfying the two-sided source equations \texttt{PjKiPj}.
Two-site support monotonicity transports the positive commuting bond from the
canonical support to the printed projection $P_x$. The local compatibility
obligation among \texttt{AppUkU=rl}, \texttt{Appetakhetc}, and
\texttt{PjKiPj} is therefore complete for both one sector and a pairwise
orthogonal family of sectors.

The source-facing theorem retains the five source identities
\begin{align}
    &\texttt{AppKxKy=0},
    \qquad
    \texttt{AppUkU=rl},
    \qquad
    \texttt{Appetakhetc},
    \label{eq:cpsv16_gsnnch_sector_decomposition_source_identity_package_first}\\
    &\texttt{Apptralktrrk},
    \qquad
    \texttt{AppPsiPhi}.
    \label{eq:cpsv16_gsnnch_sector_decomposition_source_identity_package_second}
\end{align}
These identifiers refer to the five-identity package used in Appendix~C.2 of
Ref.~\cite{Cirac2016MPDO_arXiv}. The theorem also states the standing Case~II
assumptions: the full tensor generates positive semidefinite operators, and
the simultaneous one-site tuples of its BNT representatives span the product
of their virtual matrix algebras. Every BNT bond dimension is positive. The
bare sector-decomposition type does not store this positivity, which is
necessary because algebraic injectivity is vacuous in dimension zero.

The five identities in
Eqs.~(\ref{eq:cpsv16_gsnnch_sector_decomposition_source_identity_package_first})
and~(\ref{eq:cpsv16_gsnnch_sector_decomposition_source_identity_package_second})
do not imply common copy weights. Appendix~C.2 of
Ref.~\cite{Cirac2016MPDO_arXiv} invokes \texttt{lemmus} to obtain this equality
from zero correlation length, but Proposition~\texttt{prop3to4} of the same
source does not assume zero correlation length. The formal proof therefore
does not use \texttt{lemmus}.

At each fixed length $N\geq2$, sitewise compression by the pairwise orthogonal
physical supports isolates
\begin{align}
    q_s(N)\rho^{(N)}(A_s),
    \label{eq:cpsv16_gsnnch_sector_decomposition_compressed_sector_term}\\
    q_s(N)
        &=\sum_q\mu_{s,q}^N.
    \label{eq:cpsv16_gsnnch_sector_decomposition_sector_power_sum_coefficient}
\end{align}
Global positivity makes the compressed operator in
Eq.~(\ref{eq:cpsv16_gsnnch_sector_decomposition_compressed_sector_term})
positive. The sector operator is positive and nonzero, so uniqueness of its
positive scalar phase implies that $q_s(N)$ is a nonnegative real number.
Define the bond-rescaling factor by
\begin{align}
    r_s(N)
        &=\left(\frac{q_s(N)}{n_s}\right)^{1/N}.
    \label{eq:cpsv16_gsnnch_sector_decomposition_positive_root_bond_rescaling}
\end{align}
Rescaling the supported commuting bond by $r_s(N)$ preserves the natural copy
multiplicity $n_s$ and reproduces the exact unnormalized sector term.
Definition~4.8 of Ref.~\cite{Cirac2016MPDO_arXiv} chooses its data separately
at every length, so this construction proves the proposition without copy
independence.

The local Markov labels in Appendix~C.2 of
Ref.~\cite{Cirac2016MPDO_arXiv} belong within one injective BNT component.
They are not the outer labels $x$ of Definition~4.8 and must not replace those
outer sectors.

\section{Classification}

The source GSNNCH data and predicate are formalized. The represented sector
sum is positive semidefinite and translation invariant, and the translates
of each sector bond commute pairwise. The single-bond presentation gives the
one-sector case.

The exact unstarred BNT layer identity in
Eq.~(\ref{eq:cpsv16_gsnnch_sector_decomposition_unstarred_layer_annihilation})
is formalized. Under the standing biCF injectivity hypothesis and the
per-sector MPDO positivity supplied by the positive neighboring-operator
factorization, it gives pairwise orthogonal one-site projections supporting
every BNT tensor on both sides. This theorem does not assert that the
projections sum to the identity. A square-zero family shows that such a
conclusion is false for arbitrary MPO tensors without these hypotheses.

The construction from an orthogonal commuting sector family to the GSNNCH
form is available with the natural BNT multiplicities. The canonical
restriction, its exact sectorwise factorization, and its positive neighboring
operators follow from \texttt{AppUkU=rl} and \texttt{Appetakhetc}; inactive
factor directions are removed without a minimality assumption. Together with
the two-sided supports obtained from \texttt{AppKxKy=0}, this gives supported
positive commuting bonds whose periodic products are exact on the printed
projections.

Fixed-length orthogonal compression proves that the actual BNT power-sum
coefficients are nonnegative real numbers. Positive-root bond rescaling then
absorbs each coefficient while preserving the natural BNT copy number. The
source-facing theorem retains \texttt{Apptralktrrk} and \texttt{AppPsiPhi} in
its statement.

Proposition~\texttt{prop3to4} of Ref.~\cite{Cirac2016MPDO_arXiv} is therefore
complete under its printed five identities together with the standing Case~II
assumptions: MPDO positivity, simultaneous one-site spanning of the product of
the BNT virtual matrix algebras, and positive BNT bond dimensions. This
chainwise argument repairs the proposition. It does not validate the printed
invocation of Lemma~\texttt{lemmus} and does not prove copy independence.

\bibliographystyle{alpha}
\bibliography{references}

\end{document}
